\section{Rational instruments stable under external references}\label{sec:reference66}
A rounding operation on density matrices need not be affine. In
particular, the map of Lemma~\ref{lem:positiveround65} cannot be tensored
with the identity on an unknown reference system. We therefore round a
\emph{description of the instrument}, once and independently of its
input, rather than round the state on which it acts. The resulting
maps are completely positive and jointly trace preserving. All
constants below are explicit in the input and output dimensions and
the number of outcomes.

\subsection{An integral repair of Choi data}
Write $M_d$ for the complex matrices of order $d$. An instrument from
$M_d$ to $M_n$ with $m$ outcomes is a family of completely positive
maps $\Phi_y$ whose sum is trace preserving. We use the input-first,
\emph{unnormalized} Choi convention
\begin{equation}\label{eq:choiconvention66}
 J_y=\sum_{i,j=1}^d |i\rangle\langle j|\otimes
                     \Phi_y(|i\rangle\langle j|),\qquad
 J_y\ge0,\qquad \sum_{y=1}^m\operatorname{Tr}_{n}J_y=I_d.
\end{equation}
Here $\operatorname{Tr}_{n}$ traces the output factor; in particular
$\sum_y\tr J_y=d$, not one. These are the classical Choi criteria
\cite[Theorems~2.22 and~2.26]{Watrous65}. The channel with the outcome
retained is $\mathcal F(X)=\bigoplus_y\Phi_y(X)$. We write $\|\cdot\|_\diamond$
for the completely bounded trace norm.

\begin{lemma}[Choi trace norm and reference stability]\label{lem:choidiamond66}
For any linear map $\Lambda:M_d\to M_n$ with unnormalized Choi matrix
$J(\Lambda)$,
\[
 d^{-1}\|J(\Lambda)\|_1\le\|\Lambda\|_\diamond
                                      \le\|J(\Lambda)\|_1.
\]
Consequently the diamond distance of two instruments, with the outcome
retained, is at most the sum of the trace norms of their Choi-block
differences.
\end{lemma}
\begin{proof}
Applying $\operatorname{id}_d\otimes\Lambda$ to the normalized
maximally entangled state gives the lower bound. For the upper bound,
let $u,v$ be unit vectors in $\C^r\otimes\C^d$. With
$\Omega=\sum_i e_i\otimes e_i$, write
$u=(A\otimes I_d)\Omega$ and $v=(B\otimes I_d)\Omega$.
The matrices $A,B:\C^d\to\C^r$ have Hilbert--Schmidt norm one and
therefore operator norm at most one. Thus
\[
 (\operatorname{id}_r\otimes\Lambda)(|u\rangle\langle v|)
       =(A\otimes I_n)J(\Lambda)(B^*\otimes I_n)
\]
has trace norm at most $\|J(\Lambda)\|_1$. The trace-norm unit ball
is the convex hull of rank-one partial isometries (including zero as a
convex combination of opposite such matrices). A singular-value
decomposition of an arbitrary input matrix therefore proves the same induced
trace-norm bound, independently of $r$. Taking the supremum over $r$
proves the assertion. For an instrument the Choi matrix is, up to a
fixed permutation of factors, the direct sum of its blocks.
\end{proof}

Let $B$ be a positive integer. The supplied approximations $T_y$ are
Hermitian matrices with entries in $B^{-1}\Z[i]$. Assume that every
independent real diagonal coordinate, and every independent real and
imaginary off-diagonal coordinate, differs from that of $J_y$ by at
most $B^{-1}$. The $T_y$ need not be positive or satisfy the marginal
constraint. Put
\begin{align}
 R&=I_d-\sum_y\operatorname{Tr}_{n}T_y,&
 Q_y&=T_y+\mathbf1_{\{y=1\}}R\otimes|e_1\rangle\langle e_1|,
                                                    \label{eq:choirepair66}\\
 c&=2nd(1+mn),&D&=B+mnc,\qquad
 \widetilde J_y=\frac{BQ_y+cI_{dn}}{D}.\label{eq:choigrid66}
\end{align}
The choice of the first outcome and first output coordinate is merely
a fixed convention. Any fixed outcome and output basis coordinate give
the same legality and uniform error bound. If specified zero outcomes
must remain exactly zero, remove them before applying the theorem and
restore them afterwards; $m$ then counts only the repaired outcomes.

\begin{theorem}[Reference-stable rational instrument rounding]\label{thm:choiround66}
For every $d,n,m,B\ge1$, the matrices in \eqref{eq:choigrid66} are
Choi matrices of an instrument. Their numerators are positive semidefinite matrices with Gaussian-integer entries, their common denominator is $D$, and
\begin{equation}\label{eq:choierror66}
 \|\mathcal F-\widetilde{\mathcal F}\|_\diamond
 \le\sum_y\|J_y-\widetilde J_y\|_1
 \le\frac{3mndc}{B+mnc}\le\frac{K(d,n,m)}B,
 \qquad K(d,n,m)=6mn^2d^2(1+mn).
\end{equation}
The construction uses integer additions, multiplication by specified
integers, and the supplied grid entries. It uses no spectral
factorization, positivity projection, or rational Kraus
factorization. Its complete output description has
\[
 O\bigl(m d^2n^2\log(B+mnc+1)\bigr)
\]
bits, with an absolute implicit constant for the fixed encoding.
\end{theorem}
\begin{proof}
Set $E_y=T_y-J_y$. Counting the real coordinates gives
$\|E_y\|_{\rm op}\le\|E_y\|_F\le2nd/B$.
For a Hermitian $E$, the inequalities
$-\|E\|_{\rm op}I\le E\le\|E\|_{\rm op}I$ imply
$\|\operatorname{Tr}_{n}E\|_{\rm op}\le n\|E\|_{\rm op}$.
It follows that
\[
 \|R\|_{\rm op}\le2mn^2d/B,\qquad
             \|Q_y-J_y\|_{\rm op}\le c/B.
\]
Since $J_y\ge0$, the numerator $BQ_y+cI_{dn}$ is positive.
All its entries are Gaussian integers. The partial trace constraint is
exact, because
\[
 \sum_y\operatorname{Tr}_{n}(BQ_y+cI_{dn})
                       =(B+mnc)I_d.
\]
The Choi criteria therefore give a completely positive instrument,
including when some of the original outcomes are zero.

Let $F_y=Q_y-J_y$. The identity
\[
 D(\widetilde J_y-J_y)=BF_y+cI_{dn}-mncJ_y
\]
and $\|F_y\|_1\le dn c/B$ give
\[
 D\sum_y\|\widetilde J_y-J_y\|_1
 \le mndc+mndc+mnc\sum_y\tr J_y=3mndc.
\]
Lemma~\ref{lem:choidiamond66} proves the norm assertion. Finally each
diagonal entry of a positive numerator is at most $D$, by the summed
partial trace constraint. Positivity bounds every off-diagonal modulus
by the geometric mean of two diagonal entries, hence also by $D$.
Encoding $m(dn)^2$ entries and the common denominator proves the
bit bound. All intermediate entries have polynomially bounded bit
length in the stated parameters and the supplied data.
\end{proof}

For exact rational input, directed truncation supplies $T_y$ by integer
division. For uncertain real input, the hypothesis is a certified error
bound on the \emph{final grid coordinates}; an uncertified floating
estimate is not a premise of the theorem. Although $B$ may be a power
of two, $D=B+mnc$ generally is not. The output is rational, not in
general dyadic. The construction may replace a zero outcome by a
small nonzero one; its mass is charged in \eqref{eq:choierror66}.
Known zero outcomes can instead be removed before rounding and restored
as zero blocks afterwards, with $m$ counting the remaining outcomes.

\begin{corollary}[A finite-description instrument net]\label{cor:instrumentnet66}
For fixed $d,n,m$ and $0<\delta\le1$, the preceding construction yields
a finite diamond-$\delta$ net of legal instruments. Each member has
an integer description of length
\[
 O\bigl(m d^2n^2\log(2+K(d,n,m)/\delta+mnc)\bigr).
\]
This is a deterministic description net, not a statistically learned net.
A particular member is computable in polynomial bit time from rational
input data and the grid precision; enumeration of the entire net is
not required.
\end{corollary}
\begin{proof}
Every Choi diagonal is at most one by \eqref{eq:choiconvention66}, and
positivity bounds every entry modulus by one. There are finitely many
possible directed grid truncations at $B=\lceil K/\delta\rceil$.
Apply the theorem to those arising from valid instruments. Computation
of one truncation and its repair uses a polynomial number of elementary
integer operations on polynomial-length operands.
\end{proof}
The CP/TP constraints and adding a scalar identity buffer are not new
principles. Projected least-squares process tomography already repairs
Choi estimates and uses identity mixing to finish a TP-preserving
positivity correction \cite[Section~4]{PLS66}. Here the particular
integer formula supplies an explicit denominator and a deterministic
diamond-error budget from coordinate errors. No nearest-projection,
statistical sample-complexity, or rank-optimality claim is made.

\subsection{Adaptive use and stopped observations}
A tester may retain an arbitrary finite-dimensional quantum memory,
start with a state entangled across that memory and the instrument
input, and insert arbitrary channels and finite-outcome measurements between uses.
The command port is classical: its value may be selected from the
public history and measured tester memory. Coherent access to a
superposition of commands is not an interface assumed here. The
\emph{same} tester is used in the two compared experiments.

\begin{theorem}[Adaptive reference-stable precision budget]\label{thm:adaptivediamond66}
Suppose that at use $t$, uniformly over permitted commands and public
histories,
$\|\mathcal F_{t,a}-\widetilde{\mathcal F}_{t,a}\|_\diamond\le\delta_t$.
For every common tester with a classical command port (not coherent
superpositions of commands), every initial state, and every public stopping
rule $\tau\le N$, the final states, with the stopped history and tester
memory retained, satisfy
\begin{equation}\label{eq:testererror66}
 \|\Omega_\tau-\widetilde\Omega_\tau\|_1
                  \le\sum_{t=1}^N\delta_t.
\end{equation}
The bound is independent of the tester-memory dimension. In particular,
rounding each command by Theorem~\ref{thm:choiround66} with
$B\ge N K(d,n,m)2^L$ gives stopped joint-state error at most $2^{-L}$
and transcript total-variation error at most $2^{-L-1}$.
\end{theorem}
\begin{proof}
Replace the instruments one use at a time. Just before the replaced
use, both members of a neighboring hybrid pair have the same joint
state. Decomposing a classical command/history register into blocks,
the diamond bound applies to each positive block and is multiplied by
its trace. These traces sum to one. The error introduced at this use
is therefore at most $\delta_t$, including any quantum reference.
All subsequent operations in the hybrid pair are the same channels;
complete trace-norm contraction bounds their final difference by the
same quantity. Summing the $N$ hybrid differences proves
\eqref{eq:testererror66}. For a stopping rule retain an absorbing halt
flag, use the identity on halted blocks, and finish all experiments at
use $N$. A public history-dependent stopping rule is then part of the
same tester. The same argument applies, without conditioning on a
possibly rare stopping history. Taking traces and dividing by two
gives the total-variation assertion.
\end{proof}
The hybrid principle is classical; see
\cite[Section~3.3, equation~(3.306)]{Watrous65} and the operational
strategy norms of \cite{Gutoski66}. The conclusion here is a uniform
finite-description approximation by \emph{genuine instruments}. It is
not a classical device that accepts an unknown physical entangled
state and prepares its output. A numerical simulation of a specified
joint state must charge the joint matrix dimension. Nor does the
nonlinear state repair in the preceding section become a channel by
this theorem.

\begin{proposition}[Posterior control at a stopping time]\label{prop:posterior66}
Let $A_h=p_h\rho_h$ and $\widehat A_h=q_h\widehat\rho_h$ be the
positive blocks at a common bounded stopping time, with both total
traces one. At zero simulated mass assign any density matrix to
$\widehat\rho_h$. If
$\sum_h\|A_h-\widehat A_h\|_1\le\delta$, then
\begin{equation}\label{eq:posterior66}
 \sum_h p_h\|\rho_h-\widehat\rho_h\|_1\le2\delta,\qquad
 \Pr_p\{\|\rho_h-\widehat\rho_h\|_1>\eta\}
                              \le\min\{1,2\delta/\eta\}.
\end{equation}
The tail estimate holds for every $\eta>0$. For any event of true probability $\alpha>0$, the two pooled
conditional states have distance at most $\min\{2,2\delta/\alpha\}$,
with an arbitrary conditional state when the simulated event has
probability zero.
\end{proposition}
\begin{proof}
Writing $e_h=\|A_h-\widehat A_h\|_1$, trace contraction gives
$|p_h-q_h|\le e_h$. For positive $p_h$,
\[
 p_h\|\rho_h-\widehat\rho_h\|_1
 \le e_h+|p_h-q_h|\le2e_h.
\]
The inequality remains valid if $q_h=0$. Sum and apply Markov's
inequality. Applying the same calculation to the sums over an event
gives the final assertion; two density matrices have distance at most
two. Zero true mass makes no contribution.
\end{proof}
The proposition also applies to the numerical history metric of
Theorem~\ref{thm:instrumentstream65}. For its stopped version, freeze
every halted block and apply the same contraction and mass-weighted
repair argument only to live blocks. At each use their total mass is
at most one, so the bound remains $(N+1)\delta_B$. Thus the theorem
provides a high-probability posterior guarantee at a chosen stopping
time without a probability floor or a union bound over all times.
It does not assert simultaneous accuracy of every rare normalized
posterior. This distinction is relevant when comparing finite-bit
history approximation with numerical quantum filtering
\cite[Section~II, equations~(10)--(13)]{RouchonRalph66}.
